\documentclass[12pt,a4paper]{article}
% ---------------- Paquetes ----------------
\usepackage[spanish]{babel}
\usepackage[utf8]{inputenc}
\usepackage{amsmath}
\usepackage{geometry}
\usepackage{hyperref}
\usepackage{listings}
\usepackage{xcolor}
\usepackage{graphicx}
\usepackage{float}

% Márgenes
\geometry{margin=2.5cm}

% Definición del lenguaje Dart (listings no lo incluye)
\lstdefinelanguage{Dart}{
  morekeywords={abstract,as,async,await,break,case,catch,class,const,continue,
    default,do,else,enum,extends,false,final,finally,for,get,if,import,in,is,
    late,new,null,override,required,return,set,static,super,switch,this,throw,
    true,try,var,void,while,with,int,double,String,bool,Widget,State},
  sensitive=true,
  morecomment=[l]{//},
  morecomment=[s]{/*}{*/},
  morestring=[b]',
  morestring=[b]"
}

% Configuración para código
\lstset{
  language=Dart,
  backgroundcolor=\color{gray!10},
  basicstyle=\ttfamily\footnotesize,
  keywordstyle=\color{blue},
  commentstyle=\color{green!60!black},
  stringstyle=\color{red},
  showstringspaces=false,
  frame=single,
  breaklines=true,
  literate={á}{{\'a}}1 {é}{{\'e}}1 {í}{{\'i}}1 {ó}{{\'o}}1 {ú}{{\'u}}1
           {Á}{{\'A}}1 {É}{{\'E}}1 {Í}{{\'I}}1 {Ó}{{\'O}}1 {Ú}{{\'U}}1
           {ñ}{{\~n}}1 {Ñ}{{\~N}}1 {¡}{{!`}}1 {¿}{{?`}}1
           {×}{{$\times$}}1 {÷}{{$\div$}}1
}

% Captura de pantalla (si el archivo no existe, muestra un recuadro)
\newcommand{\captura}[2]{%
  \begin{figure}[H]
    \centering
    \IfFileExists{capturas/#1}%
      {\includegraphics[height=0.42\textheight]{capturas/#1}}%
      {\fbox{\parbox{0.5\linewidth}{\centering\vspace{3cm}Insertar captura:\\\texttt{capturas/#1}\vspace{3cm}}}}
    \caption{#2}
  \end{figure}}

% ---------------- Portada ----------------
\newcommand{\data}{2026} % fecha
\newcommand{\subtitle}{Practica 03}
\newcommand{\titulo}{DESARROLLO DE SOFTWARE II} 
\newcommand{\departmento}{Universidad Nacional de San Antonio Abad del Cusco}
\newcommand{\curso}{Escuela Profesional de Ingeniería Informática y de Sistemas}

\begin{document}

% ----------- Portada -----------
\begin{titlepage}
\begin{center}\leavevmode
    \normalfont
    \textsc{\Large \departmento}\\[1 cm]
    {\large \curso}
    \vskip 1cm
    \includegraphics[width=0.40\columnwidth]{unsaac.png}
    \includegraphics[width=0.40\columnwidth]{logo.png}
    \rule{\linewidth}{0.2 mm}
    {\huge \bfseries \titulo \par}{}
    \vskip 1cm
    {\Large \bfseries \subtitle \par}
           
    \rule{\linewidth}{0.2 mm}\\[1.0 cm]
     
    \begin{minipage}[t]{0.45\textwidth}
    	\begin{flushleft} \large
    		\emph{Alumno:}\\
    		Ibeth Janela del Pilar Cusi Diaz\\
            \emph{Docente:}\\
            Hans Harley Ccacyahuillca Bejar
    	\end{flushleft}
    \end{minipage}
    \vfill
    {\normalsize \data\par}
\end{center}
\end{titlepage}

% ----------- Contenido -----------
\tableofcontents
\newpage

\section{Objetivos}
\begin{itemize}
  \item Conocer el entorno de desarrollo Android Studio y su integración con Flutter.
  \item Instalar y configurar el SDK de Flutter y las dependencias necesarias.
  \item Comprender la estructura de un proyecto Flutter y el rol del \textit{widget tree}.
  \item Crear y ejecutar la primera aplicación Flutter en un emulador Android.
  \item Explorar los widgets fundamentales: \texttt{Text}, \texttt{Container}, \texttt{Row}, \texttt{Column} e \texttt{Image}.
\end{itemize}

\section{Marco teórico}
Flutter es un framework de código abierto creado por Google para construir aplicaciones compiladas para móvil, web y escritorio a partir de una única base de código escrita en Dart. No utiliza componentes nativos del sistema operativo, sino que dibuja su propia interfaz con el motor gráfico Skia/Impeller. En Flutter todo es un \textit{widget}, y los widgets se organizan en un árbol jerárquico (\textit{widget tree}). Un \texttt{StatelessWidget} es inmutable, mientras que un \texttt{StatefulWidget} mantiene un estado que se actualiza con \texttt{setState()}.

\section{Entorno y configuración}
Se instaló el SDK de Flutter (canal \textit{stable}), se añadió \texttt{flutter\textbackslash bin} al \texttt{PATH}, se instaló el plugin de Flutter en Android Studio (que incluye Dart) y se creó un emulador Pixel con API 33. La instalación se verificó con \texttt{flutter doctor}.

\captura{flutterdoctor.png}{Resultado del comando \texttt{flutter doctor}.}

\section{Organización de la aplicación}
Los tres ejercicios y el ejercicio propuesto se integraron en \textbf{una sola aplicación}. Cada ejercicio está en su propio archivo (para verlos por separado en el repositorio Git) y \texttt{home\_screen.dart} los une mediante una \texttt{NavigationBar} y un \texttt{IndexedStack}, que conserva el estado de cada pestaña.

\begin{lstlisting}[language={},caption={Estructura del proyecto.}]
lib/
|-- main.dart
|-- logic/calculadora_logica.dart
`-- screens/
    |-- home_screen.dart
    |-- ejercicio1_hola.dart
    |-- ejercicio2_perfil.dart
    |-- ejercicio3_contador.dart
    `-- ejercicio4_calculadora.dart
\end{lstlisting}

\lstinputlisting[caption={\texttt{main.dart}: punto de entrada.}]{../lib/main.dart}
\lstinputlisting[caption={\texttt{home\_screen.dart}: unión de los ejercicios.}]{../lib/screens/home_screen.dart}
\captura{home.png}{Aplicación con la barra de navegación que une los ejercicios.}

\section{Ejercicio 1: Primer proyecto Flutter}
Se creó el proyecto \texttt{mi\_primera\_app} y se construyó una pantalla con \texttt{Scaffold}, \texttt{AppBar}, \texttt{Container}, \texttt{Column}, \texttt{Row}, \texttt{Text} e \texttt{Image}. Si no existe el recurso \texttt{assets/images/logo.png}, \texttt{errorBuilder} muestra el \texttt{FlutterLogo}.

\lstinputlisting[caption={Ejercicio 1.}]{../lib/screens/ejercicio1_hola.dart}
\captura{ejercicio1.png}{Ejecución del Ejercicio 1 en el emulador.}

\section{Ejercicio 2: Tarjeta de perfil}
Se utilizaron \texttt{Card}, \texttt{CircleAvatar}, \texttt{Divider}, \texttt{Row} con \texttt{spaceEvenly} y un widget auxiliar reutilizable \texttt{\_Estadistica}. Se agregó un \texttt{SingleChildScrollView} para evitar desbordes en pantallas pequeñas.

\lstinputlisting[caption={Ejercicio 2.}]{../lib/screens/ejercicio2_perfil.dart}
\captura{ejercicio2.png}{Ejecución del Ejercicio 2 en el emulador.}

\section{Ejercicio 3: Contador interactivo}
Se implementó un \texttt{StatefulWidget} con tres botones (incrementar, decrementar y reiniciar). El color del número cambia según su valor: gris en 0, azul de 1 a 4 y verde desde 5.

\lstinputlisting[caption={Ejercicio 3.}]{../lib/screens/ejercicio3_contador.dart}
\captura{ejercicio3.png}{Ejecución del Ejercicio 3 en el emulador.}

\section{Ejercicio propuesto: Calculadora básica}
\subsection{Diseño}
La lógica se separó de la interfaz: \texttt{CalculadoraLogica} (Dart puro) gestiona los dígitos, el punto decimal, las cuatro operaciones y el manejo de errores (división entre cero). Las operaciones encadenadas se evalúan de izquierda a derecha, como una calculadora básica. La pantalla (\texttt{StatefulWidget}) dibuja el \textit{display} y un \texttt{GridView} de 4$\times$4 botones, además del botón \texttt{C}.

El color del \textit{display} depende del valor mostrado:
\begin{itemize}
  \item Positivo: azul.
  \item Negativo: rojo.
  \item Cero: gris.
\end{itemize}

\lstinputlisting[caption={Lógica de la calculadora.}]{../lib/logic/calculadora_logica.dart}
\lstinputlisting[caption={Interfaz de la calculadora.}]{../lib/screens/ejercicio4_calculadora.dart}

\captura{calc1.png}{Resultado positivo (azul).}
\captura{calc2.png}{Resultado negativo (rojo).}
\captura{calc3.png}{Display en cero (gris).}

\subsection{Pruebas}
Se escribieron pruebas unitarias para la suma, decimales, resultados negativos, operaciones encadenadas, división entre cero y limpieza, ejecutadas con \texttt{flutter test}.

\lstinputlisting[caption={Pruebas unitarias.}]{../test/widget_test.dart}

\section{Cuestionario}
\subsection*{1. ¿Diferencia entre StatelessWidget y StatefulWidget?}
Un \texttt{StatelessWidget} es inmutable: su apariencia depende solo de los datos con que fue creado y no cambia por sí mismo. Un \texttt{StatefulWidget} tiene un objeto \texttt{State} asociado, que persiste entre reconstrucciones y se modifica con \texttt{setState()}. Ejemplos: un título o una etiqueta (\texttt{Text}, \texttt{\_Estadistica}) es \textit{stateless}; el contador o la calculadora de esta práctica son \textit{stateful}.

\subsection*{2. ¿Por qué Flutter no usa componentes nativos?}
Porque dibuja directamente cada píxel con su propio motor gráfico (Skia/Impeller) sobre un lienzo. Esto da una apariencia idéntica en todas las plataformas, evita puentes entre Dart y los componentes nativos, ofrece control total sobre la interfaz y permite alto rendimiento (60 o 120 fps).

\subsection*{3. ¿Qué hace \texttt{flutter doctor}?}
Revisa que el entorno esté bien configurado. Sus secciones principales son:
\begin{itemize}
  \item \textbf{Flutter}: versión y canal del SDK.
  \item \textbf{Android toolchain}: SDK de Android, \textit{build-tools} y licencias aceptadas.
  \item \textbf{Android Studio}: IDE instalado y plugins de Flutter y Dart.
  \item \textbf{Connected device}: emuladores o dispositivos físicos disponibles.
\end{itemize}
Cada elemento correcto aparece con un check verde; los problemas se muestran con la acción sugerida.

\subsection*{4. ¿Qué sucede al llamar \texttt{setState()}?}
\texttt{setState()} ejecuta la función recibida para modificar el estado y marca el \texttt{State} (su \texttt{Element}) como \textit{dirty}, programándolo para el siguiente \textit{frame}. En ese frame Flutter vuelve a ejecutar \texttt{build()} de ese widget y obtiene un subárbol nuevo de widgets. Luego compara cada widget nuevo con el anterior (mismo tipo y misma \texttt{key}): si coinciden, el \texttt{Element} se reutiliza y solo se actualizan los \texttt{RenderObject} cuyas propiedades cambiaron; si no, se crea un elemento nuevo. Así solo se repinta lo necesario.

\subsection*{5. Hot Reload vs. Hot Restart}
\textbf{Hot Reload} inyecta el código modificado en la máquina virtual de Dart en ejecución y reconstruye el árbol de widgets, conservando el estado (por ejemplo, el valor del contador). \textbf{Hot Restart} reinicia la aplicación desde \texttt{main()}, reinicia el estado y es necesario al cambiar \texttt{main()}, \texttt{initState()}, variables globales o estáticas y enumeraciones.

\section{Conclusiones}
\begin{itemize}
  \item Se configuró el entorno Flutter en Android Studio y se ejecutó la aplicación en un emulador Android.
  \item El uso de widgets de layout (\texttt{Row}, \texttt{Column}, \texttt{Card}, \texttt{GridView}) permite componer interfaces de forma declarativa.
  \item Separar la lógica (\texttt{CalculadoraLogica}) de la interfaz facilita las pruebas y el mantenimiento.
  \item Organizar cada ejercicio en un archivo y unirlos en una sola app mejora la claridad del repositorio.
\end{itemize}

\section{Repositorio}
Código fuente: \url{https://github.com/USUARIO/REPOSITORIO} \quad (reemplazar por el enlace real).

\begin{thebibliography}{9}
\bibitem{flutter} Google LLC. \textit{Flutter Documentation --- Get Started}. \url{https://docs.flutter.dev/get-started/install}
\bibitem{dart} Google LLC. \textit{Dart Language Tour}. \url{https://dart.dev/language}
\bibitem{nahavandipoor} V. Nahavandipoor. \textit{Flutter \& Dart Cookbook}. O'Reilly Media, 2023.
\bibitem{windmill} E. Windmill. \textit{Flutter in Action}. 2nd ed. Manning Publications, 2023.
\bibitem{napoli} M. Napoli. \textit{Beginning Flutter: A Hands-On Guide to App Development}. Wiley, 2020.
\end{thebibliography}

\end{document}
